\section{Count Nodes in a Complete Binary Tree}
\label{sec:trees-count-complete}

\problemstatement{Given the root of a \emph{complete} binary tree (every
level full except possibly the last, which is filled left to right), count
its nodes in better than $O(n)$ time.}

\begin{patternbox}
A plain traversal counts in $O(n)$, but the word \emph{complete} is an
invitation to do better. The key fact: if a subtree's leftmost depth equals
its rightmost depth, it is a \textbf{perfect} tree and its node count is
$2^{\text{height}}-1$ by formula -- no traversal needed. Only along the
one ``ragged'' edge do we recurse, giving $O(\log^2 n)$.
\end{patternbox}

\subsection*{Understanding the problem carefully}

Walk down the far left of a subtree to get its left height; walk down the
far right to get its right height. If they are equal, every level is full,
so the subtree is perfect and has exactly $2^h-1$ nodes -- return that
instantly. If they differ, the last level is only partially filled, so
recurse into both children and add $1$ for the current node. Crucially, at
most one child per level is ragged; the other is perfect and resolved by
formula, so the recursion depth is $O(\log n)$ and each level costs
$O(\log n)$ to measure the two heights.

\begin{ideabox}[Perfect Subtrees Close by Formula]
\textsc{Count}(node): if null return $0$. Measure left height $l$
(all-left walk) and right height $r$ (all-right walk). If $l=r$, return
$2^{l}-1$. Otherwise return $1 + \textsc{Count}(\text{left}) +
\textsc{Count}(\text{right})$.
\end{ideabox}

\subsection*{Algorithm}

\begin{enumerate}
  \item \textsc{leftHeight}(node): count nodes going only left until null.
  \item \textsc{rightHeight}(node): count nodes going only right until
        null.
  \item \textsc{Count}(node): if null return $0$; if leftHeight $=$
        rightHeight $=h$ return $2^h-1$; else return $1+$
        \textsc{Count}(left)$+$\textsc{Count}(right).
\end{enumerate}

\subsection*{Dry run}

\dryrun{Complete tree: root $1$; children $2,3$; $2$'s children $4,5$;
$3$'s left child $6$.}

\begin{itemize}
  \item At $1$: left height $=3$ ($1\!\to\!2\!\to\!4$), right height $=2$
        ($1\!\to\!3$). Unequal $\Rightarrow$ recurse.
  \item Left subtree at $2$: left height $=$ right height $=2$; perfect,
        $2^2-1=3$ nodes.
  \item Right subtree at $3$: left height $=2$, right height $=1$; recurse
        $\to 1+\textsc{Count}(6)+0 = 1+1 = 2$.
  \item Total $=1+3+2=6$.
\end{itemize}

\subsection*{C++ implementation}

\begin{lstlisting}
#include <bits/stdc++.h>
using namespace std;

struct TreeNode {
    int val;
    TreeNode* left;
    TreeNode* right;
    TreeNode(int v) : val(v), left(nullptr), right(nullptr) {}
};

int leftHeight(TreeNode* n) {
    int h = 0;
    while (n) { h++; n = n->left; }
    return h;
}
int rightHeight(TreeNode* n) {
    int h = 0;
    while (n) { h++; n = n->right; }
    return h;
}

int countNodes(TreeNode* root) {
    if (root == nullptr) return 0;
    int lh = leftHeight(root), rh = rightHeight(root);
    if (lh == rh) return (1 << lh) - 1;          // perfect subtree
    return 1 + countNodes(root->left) + countNodes(root->right);
}

int main() {
    TreeNode* root = new TreeNode(1);
    root->left = new TreeNode(2);
    root->right = new TreeNode(3);
    root->left->left = new TreeNode(4);
    root->left->right = new TreeNode(5);
    root->right->left = new TreeNode(6);

    cout << "Node count: " << countNodes(root) << "\n";
    return 0;
}
\end{lstlisting}

\begin{complexitybox}
\textbf{Time Complexity.} $O(\log^2 n)$ -- the recursion has $O(\log n)$
depth, and each level spends $O(\log n)$ measuring the two heights.
\textbf{Space Complexity.} $O(\log n)$ for the recursion stack.
\end{complexitybox}

\begin{takeawaybox}
When a structural guarantee (``complete'') is handed to you, look for a
closed-form shortcut before defaulting to a full traversal. Resolving
perfect subtrees by the $2^h-1$ formula and recursing only along the ragged
boundary is the canonical sub-linear tree count.
\end{takeawaybox}
